% ============= Angaben zur Arbeit =============

\newcommand{\titel}{Xxxxx Xxxxx Xxxxx Xxxxx Xxxxx Xxxxx Xxxxx Xxxxx Xxxxx Xxxxx Xxxxx} % Hier den Titel der Arbeit angeben 
\newcommand{\arbeit}{MASTERARBEIT} % Hier die Art der Abschlussarbeit eingeben in Großbuchstaben 
\newcommand{\Nummer}{XX} % Nummerrierung der Arbeit (mit Berteurer zu klären)
\newcommand{\EndeMonatTxt}{Xxxxx} % Angabe des Monates der Abgabe als Wort
\newcommand{\EndeJahr}{20XX} % Angabe des Jahres der Abgabe
\newcommand{\bereuer}{Xxxxx Xxxxx} % Hier den Betreuer der Arbeit angeben
\newcommand{\pruefer}{Xxxxx Xxxxx} % Hier den Prüfer der Arbeit angeben 
\newcommand{\student}{Xxxxx Xxxxx} %Hier den Studierendenname der Arbeit angeben 
\newcommand{\profEins}{Prof. Dr.-Ing. M.Arch. Lucio Blandini} % Hier den Professor1 des Instituts angeben 
\newcommand{\profZwei}{Prof. Dr.-Ing. Dr.-Ing. E.h. Dr. h.c. Werner Sobek} % Hier den Professor2 des Instituts angeben 
\newcommand{\profDrei}{Prof. Dr.-Ing. Balthasar Nov\'{a}k} % Hier den Professor3 des Instituts angeben 
\newcommand{\institut}{Institut f{\"u}r Leichtbau Entwerfen und Konstruieren} % Hier das Institut angeben
\newcommand{\titelbild}{Images/Titelseite/Titelbild} % Hier den Pfad des Titelbildes eingeben


% ============= Pakete importieren =============

%Bei Fragen und Anregungen zu dieser Vorlage können Sie sich gerene bei Silas Kalmbach (silas.kalmbach@ilek.uni-stuttgart.de) melden.

% ============= Einseitig =============
\documentclass[a4paper,12pt]{scrreprt}
\usepackage[left= 2.5cm, right = 2cm, top=2.5cm]{geometry}

% ============= Doppelseitig =============
%\documentclass[a4paper,12pt, openright, twoside]{scrreprt}
%\usepackage[left= 2.5cm, right = 2cm, top=2.5cm]{geometry}
%% -> Anpassungen an die Kopf- und Fußzeilen beachten!


% ============= Dokumentinformationen =============

\usepackage[ 
hidelinks
]{hyperref} 
\hypersetup{ 
	pdfauthor={\student}, 
	pdftitle={\titel},
	pdfsubject={\arbeit}, 
	pdfkeywords={}
}


% ============= Farbschema =============

\usepackage{color}      %Verwendung von Farben
\usepackage{xcolor}		%Definition eigener Farben

% Uni Stuttgart Colors
\definecolor{uniSlightblue}{RGB}{0,190,255}
\definecolor{uniSdarkblue}{RGB}{0,65,145}
\definecolor{uniSdarkgrey}{RGB}{62, 68, 76}
\definecolor{uniSgreyblue}{RGB}{125,198,234}
\definecolor{uniSgrey}{RGB}{159,153,152}

% Uni Stuttgart monochrom greyscale
\definecolor{uniSblack}{RGB}{0,0,0}
\definecolor{uniSdarkgrey}{RGB}{62,68,76}
\definecolor{uniSgrey}{RGB}{159,153,152}
\definecolor{uniSlightgrey}{RGB}{185,186,187}


% ============= Standart Packages und Einstellungen =============

%\usepackage[utf8]{inputenc} %Schriftkodierung PDFLaTeX
\usepackage[utf8]{luainputenc} %Schriftkodierung LuaLaTeX
\usepackage[english, ngerman]{babel} %Sprachen
\usepackage{graphicx} %Einbinden von Bildern
\usepackage{amsfonts} % zusätzliche Schriftzeichen der American Mathematical Society
\usepackage{amsmath,amssymb} % zusätzliche Schriftzeichen der American Mathematical Society
\usepackage{typearea} %Aufteilung von Rändern und Textbereich
\usepackage{tikz} %Grafiken Zeichnen
\usetikzlibrary{shapes,arrows,positioning,calc,chains,scopes}
\usepackage{pgfplots}  %Kurven, Graphen
%\usepgfplotslibrary{dateplot, fillbetween, groupplots}
\usepackage{pdfpages} %Einfügen von PDFs
\usepackage{appendix} %Anhang
\usepackage{csquotes}
\usepackage[german]{cleveref} % Mehr Möglichkeiten zum referenzieren
\usepackage{gensymb} %Symbole: \degree \celsius \perthousand \ohm \micro
\tolerance = 400  % möglicher Abstand zwischen den Wörtern erhöhen
\usepackage[section]{placeins} %\FloatBarrier Verhindert das Abbildungen in nachfolgenden Kapitel rutschen
\makeatletter\AtBeginDocument{\expandafter\renewcommand\expandafter\subsection\expandafter{\expandafter\@fb@secFB\subsection}}\makeatother %FloatBarrier für Subsection
\usepackage{subcaption} % Bilder nebeneinander
\captionsetup[subfigure]{list=true, labelfont=bf, position=top}


% ============= Programmierung (Visualisierung) =============

\usepackage[cache=false]{minted} %Highlighting von Code, Benötigt Python und Pygments hinterlegt in den Umgebungsvariablen (pip install Pygments) ,cmd Eingabe->pdflatex -synctex=1 -interaction=nonstopmode -shell-escape %.tex
\usemintedstyle{tango, xleftmargin=20pt, linenos, fontsize=\small, breaklines=true, tabsize=4} %borland, bw, tango, frame=leftline
\usepackage{listings}
%\lstset{
%	language=python,
%	otherkeywords={1, 2, 3, 4, 5, 6, 7, 8 ,9 , 0, -, =, +, [, ], (, ), \{, \}, :, *, !},
%}

\lstdefinestyle{mystyle}{
%	backgroundcolor=\color{uniSlightgrey},   
	commentstyle=\color{uniSlightblue},
	keywordstyle=\color{uniSdarkblue},
	numberstyle=\tiny\color{uniSlightgrey},
	stringstyle=\color{uniSgreyblue},
	basicstyle=\ttfamily\small,
	xleftmargin=20pt,
%	xrightmargin=3pt,
	breakatwhitespace=false,         
	breaklines=true,                 
	captionpos=b,                    
	keepspaces=true,                 
	numbers=left,                    
	numbersep=10pt,                  
	showspaces=false,                
	showstringspaces=false,
	showtabs=false,                  
	tabsize=4
}

\lstset{style=mystyle}


% ============= Tabelleneinstellungen =============

\usepackage{multirow}						% ermoeglicht mehrere Zeilen in einer Tabellenzeile
\usepackage{multicol} 						% mehre Spalten auf eine Seite
\usepackage{longtable, tabu}				% alternatives Tabellenpackage inkl. longtabu und Farben
\usepackage{float}							% um [H] bei float Objekten verwenden zu können
\usepackage{booktabs}						% um \toprule, \midrule und \bottomrule verwenden zu können
\usepackage{tabularx}
\usepackage{csvsimple} %Excel Improtieren als CSV 
\usepackage{siunitx} %Wissenschaftliche Angabe der Einheiten \num{number}
%für Excel to Latex Add-on
\usepackage{colortbl} %Farbige Tabellen
\usepackage{rotating} %Rotieren von Tabellen
\usepackage{bigstrut} %Für Tabellen mit Vertikalen Linien


% ============= Standart Einstellungen der Pfade =============

\graphicspath{{./Images/}{./Images/Titelseite/}} %Bilderordner


% ============= Schriftarten =============
%PDF-LaTeX
%\usepackage[T1]{fontenc} %Richtige Trennung von Umlauten
%\usepackage[scale=0.87]{tgheros}
%\usepackage{fourier}

%Lua-LaTeX
\usepackage[no-math]{fontspec}

\setmainfont{UniversforUniS}[
Extension      = .ttf,
Path           = fonts/,
UprightFont    = *55Rm-Regular,
ItalicFont     = *45LtObl-Rg,
BoldFont       = *65Bd-Regular,
BoldItalicFont = *45LtObl-Rg,
]
%\setsansfont{}[]
%\setmonofont{}[]

\addtokomafont{disposition}{\rmfamily} %Überschriften selbe Schriftart wie Text

\renewcommand\familydefault\sfdefault % Formeln in Sans-Serif
\usepackage[italic, symbolgreek]{mathastext} % Formeln in Sans-Serif
\renewcommand\familydefault\rmdefault
 % Formeln in Sans-Serif

\usepackage{setspace} %Zeilenabstand
\setstretch{1.05} %Zeilenabstand

\setlength{\parindent}{0pt} %Wenn Absatzabstand, dann Einzug unnötig
\setlength{\parskip}{7pt} %Absatzabstand



% ============= Kopf- und Fußzeile >>Einseitig<< =============

\usepackage{scrlayer-scrpage}
\automark{chapter}	%Automatische Kapitelangabe in Kopfzeile
\renewcommand*{\chaptermarkformat}{} %Kapitel ohne Kapitelnummer
\clearpairofpagestyles %aktuelle Einetsllungen aus Kopf- und Fußzeile löschen
\setkomafont{pageheadfoot}{\small}		%Font für Kopfzeile
\setkomafont{pagefoot}{\small}  		%Font für Fußzeile
\setkomafont{pagenumber}{\small} 	%Font für Seitennummer
\KOMAoptions{headsepline = true} %Kopzeile Linie

\lohead[]{}
\cohead[]{}
\rohead[]{\headmark\enskip \raisebox{0.7pt}{|}\enskip \thechapter} %Darstellung Kapitel

\lofoot[]{}
\cofoot[]{}
\rofoot[\pagemark]{\pagemark} %Seitenzahl


% ============= Kopf- und Fußzeile >>Doppelseitig<< =============

%\usepackage[automark, headsepline, autooneside=false]{scrlayer-scrpage}
%\automark[section]{chapter}	%Automatische Kapitelangabe in Kopfzeile
%\renewcommand*{\chaptermarkformat}{} %Kapitel ohne Kapitelnummer
%\renewcommand*{\sectionmarkformat}{} %Kapitel ohne Kapitelnummer
%\clearpairofpagestyles %aktuelle Einetsllungen aus Kopf- und Fußzeile löschen
%\setkomafont{pageheadfoot}{\small\sffamily}		%Font für Kopfzeile
%\setkomafont{pagefoot}{\small\sffamily}  		%Font für Fußzeile
%\setkomafont{pagenumber}{\small\sffamily} 	%Font für Seitennummer
%\KOMAoptions{headsepline = true} %Kopzeile Linie
%
%
%\lehead[]{\rightmark} %Zweiseitig
%\cohead[]{}
%\rohead[]{\leftmark\enskip \raisebox{0.7pt}{|}\enskip \thechapter} %Darstellung Kapitel
%
%\lefoot[\pagemark]{\pagemark}
%\lofoot[]{}
%\cofoot[]{}
%\rofoot[\pagemark]{\pagemark} %Seitenzahl


% ============= Literatur =============

\usepackage[									
natbib=true,									% Naturwissenschaftlich
backend=biber,									% Biber statt BibTex verwenden. In TexStudio einstellen!
style=numeric,									% Stil im Verzeichnis alt. authoryear
citestyle=numeric,	                            % Stil im Fließtext
giveninits=true,										
doi=false,										% doi nicht mit angeben
isbn=false,										% isbn nicht mit angeben
sorting=none,									% Sortierung des Verzeichnis nach Autor, Titel, Jahr
% sorting=none,									% Keine Sortierung des Verzeichnisses
]{biblatex}	

% \DeclareNameAlias{default}{family-given}
\addbibresource{Literatur.bib} %Bibliographiedateien laden

\urlstyle{same}	% url Schriftart anpassen auf die des Dokuments anpassen
\usepackage{xpatch}% author in bold font in bibliography
\xpretobibmacro{author}{\mkbibbold\bgroup}{}{}
\xapptobibmacro{author}{\egroup}{}{}
\xpretobibmacro{bbx:shortauthor}{\mkbibbold\bgroup}{}{}
\xapptobibmacro{bbx:shortauthor}{\egroup}{}{}
\xpretobibmacro{bbx:editor}{\mkbibbold\bgroup}{}{}
\xapptobibmacro{bbx:editor}{\egroup}{}{}


% ============= Besondere Abkürzungen ============= 

\usepackage[figurename=Abb., % Beschriftung von Abbildung zu Abb.
tablename=Tab., % Beschriftung von Tabelle zu Tab.
labelfont=bf]{caption} %Label in Bold
\KOMAoptions{numbers = noenddot} % der Punkt hinter den Kapiteln etc wird entfernt: 1.1. wird zu 1.1



% ============= Besondere Trennungen ============= 

\hyphenation{De-zi-mal-tren-nung}


% ============= Testing ============= 

\usepackage{blindtext}	% Beispiel Texte einfügen per \blindtext	


%%%%%%%%%%%%%%%%%%%%%%%%%%%%%%%%%%%%%%%%%%%%
% ============= Dokumentbeginn =============
%%%%%%%%%%%%%%%%%%%%%%%%%%%%%%%%%%%%%%%%%%%%

\begin{document}

\pagestyle{empty} %Seiten ohne Kopf- und Fußzeile sowie Seitenzahl
\begin{titlepage}
	\newgeometry{left=2.5cm,right=2cm,top=2.5cm,bottom=1.5cm}
	%------------------------------------------------------------
	\begin{minipage}[t]{60mm}
		\begin{flushleft}
			\hspace{2.0pt}\includegraphics[width=25mm]{Images/Titelseite/ILEK-logo}
		\end{flushleft}
	\end{minipage}
	\begin{minipage}[t]{70mm}
		\begin{flushleft}
			{\fontsize{14}{20}\sffamily{\arbeit}}
			%{\large\textsf\arbeit}
		\end{flushleft}	
	\end{minipage}
	\begin{minipage}[t]{33mm}
		\begin{flushright}
			{\fontsize{14}{20}\sffamily{\Nummer}~|~\textsf{\EndeJahr}}
		\end{flushright}
	\end{minipage}
	%------------------------------------------------------------
	\begin{center} 
		\vspace{-16.0pt}\nointerlineskip\rule{\textwidth}{0.4pt}\\ 
		\vspace{2.0pt}\nointerlineskip
	\end{center}
	%------------------------------------------------------------
	
	\vspace{10mm}
	\hspace{60mm}
	\begin{minipage}[c]{105mm}
		\begin{minipage}[t][4cm][c]{\textwidth}
			\LARGE{\textsf{\titel}}
		\end{minipage}
		\hspace{0mm} 
		
		\vspace{5mm}
		\hspace{-3.0mm} 
		\begin{tabular}{p{2.5cm}l}
			\fontsize{12}{20}\textsf{Bearbeiter:}&\textsf{\student} \\
		\end{tabular}
		
		\vspace{5mm}
		\hspace{-3.0mm} 
		\begin{tabular}{p{2.5cm}l}
			\fontsize{12}{20}\textsf{Betreuer:}&\textsf{\bereuer} \\ 
		\end{tabular}
		
		\vspace{5mm}
		\hspace{-3.0mm} 
		\begin{tabular}{p{2.5cm}l}
			\fontsize{12}{20}\textsf{Prüfer:}&\textsf{\pruefer}	\\
		\end{tabular}
		
		\vspace{11mm}
		\fontsize{12}{20}\textsf{\EndeMonatTxt~\EndeJahr}
		
		
		\begin{minipage}[t][10.5cm][c]{\textwidth}
			\includegraphics[width=1.0\textwidth]{\titelbild}
		\end{minipage}
		
		
		\begin{minipage}[c]{12mm}
			\vspace{-1mm}
			\includegraphics[width=10mm]{Images/Titelseite/unistuttgart_logo_de}
		\end{minipage}
		\begin{minipage}[c]{60mm}
			\fontsize{12}{20}\textsf{Universität Stuttgart}
		\end{minipage}
		
		\fontsize{12}{20}
		\vspace{3mm}
		\textsf{\institut} \\
		\textsf{\profEins} \\
		\textsf{\profZwei} \\
		\textsf{\profDrei} 
		
	\end{minipage}
	\cleardoublepage
\end{titlepage}
%------------------------------------------------------------

\restoregeometry

\pagestyle{plain.scrheadings}
\pagenumbering{Roman} % Römische Seitenzahlen verwenden
\addcontentsline{toc}{chapter}{Aufgabenstellung}
\includepdf[pages=-]{aufgabenstellung.pdf}  %->Optional
\chapter*{Eidesstattliche Erklärung}
\addcontentsline{toc}{chapter}{Eidesstattliche Erklärung}
\label{erklaerung}
Hiermit erkläre ich, dass ich die vorliegende Arbeit selbständig verfasst habe, dass ich keine anderen als die angegebenen Quellen benutzt und alle wörtlich oder sinngemäß aus anderen Werken übernommenen Aussagen als solche gekennzeichnet habe, dass die eingereichte Arbeit weder vollständig noch in wesentlichen Teilen Gegenstand eines anderen Prüfungsverfahrens gewesen ist, dass ich die Arbeit weder vollständig noch in Teilen bereits veröffentlicht habe und dass das elektronische Exemplar mit den anderen Exemplaren übereinstimmt. \\
\\[1.5cm]
Datum:	\hrulefill\enspace Unterschrift: \hrulefill
\\[3.5cm]
\chapter*{Vorwort}
\addcontentsline{toc}{chapter}{Vorwort}
\label{danksagungen}

\begin{description}
	\item[-] Diese Vorlage dient als grober Leitfaden zu Erstellung der Abschlussarbeit. Die Formatierung ist somit nicht zwingend umzusetzen.
	\item[-] Die Formatierung des Deckblattes sollte, soweit möglich, unverändert bleiben. 
	\item[-] Von der Gliederung der Arbeit kann abgewichen werden, solang dieses sinnig begründbar ist.
\end{description}
\vspace{10mm}
Um mit \LaTeX{} zu Arbeiten, kann z.B. die Kombination folgende Programme verwendet werden.
\begin{flalign*}
&\text{1) MiKTeX:} &&\text{\url{https://miktex.org/download}}&\\
&\text{2) TeXstudio:} &&\text{\url{https://www.texstudio.org/}}
\end{flalign*}
Alternativ besteht auch die Möglichkeit Online-Dienste zu benutzen, welche mögliche Schwierigkeiten bei der Einrichtung umgehen.

\subsection*{Empfohlene Einstellungen dieser Vorlage}
Für eine problemlose Kompilierung des \LaTeX-Dokumentes ist es notwendig, einige Einstellungen in den Editor zu übernehmen.
\begin{description}
	\item[-] Als Standard Bibliographieprogramm sollte Biber ausgewählt werden
	\item[-] Als Standardcompiler ist LuaLaTeX oder PdfLaTeX zu empfehlen
\end{description} %->Optional
\chapter*{Zusammenfassung}\addcontentsline{toc}{chapter}{Zusammenfassung}
\blindtext

\blindtext


\chapter*{Abstract}
\blindtext

\blindtext
\tableofcontents %Inhaltsverzeichnis
\cleardoubleoddpage
\pagenumbering{arabic}	% Arabische Seitenzahlen verwenden
\chapter*{Bezeichnungen und Symbole}
\addcontentsline{toc}{chapter}{Bezeichnungen und Symbole}
\begin{longtable}{
		@{}
		>{\centering}p{0.15\linewidth}
		@{}
		>{\hspace*{0pt}}p{0.845\linewidth}
		@{}
	}
	\centering
	\small
	\tabularnewline
	\multicolumn{2}{@{}l}{\textsf{\textbf{Akronyme}}} \\
	APDL & Ansys Parametric Design Language \\
	DMS & Dehnungsmessstreifen \\
	FEM & Finite-Elemente-Methode \\
	
	\tabularnewline
	\multicolumn{2}{@{}l}{\textsf{\textbf{Lateinische Buchstaben}}} \\
	\textit{a} & Erster Eintrag \\
	\textit{b} & Zweiter Eintrag \\
	
	\tabularnewline
	\multicolumn{2}{@{}l}{\textsf{\textbf{Griechische Buchstaben}}} \\
	$\alpha$ & Kontinuierlicher Temperaturabminderungsfaktor \\
	$\varepsilon$ & Dehnung \\
	$\varepsilon_\mathrm{b}$ & Rechnerische Dehnung im Vierpunktbiegeversuch \\
	
	\tabularnewline
	\multicolumn{2}{@{}l}{\textsf{\textbf{Indizes}}} \\
	aktiv & Wert im aktiven Zustand \\
	min & Minimalwert \\
	
\end{longtable}
\cleardoubleoddpage
\pagestyle{scrheadings}	% pagestyle für gesamtes Dokument aktivieren


% ============= Kapitel =============

%\chapter{Einleitung}

\label{sec:einleitung}

\blindtext
\cite{test}

\blindtext
\blindtext
\cite{kroll_computational_2016, yang_application_2003}

\blindtext
\cite{yang_application_2003}

\chapter{Float Objekte}

1 m : Leerzeichen \\
1\,m : schmales Leerzeichen \\

\begin{description}
	\item[h:] an der Stelle, an der es in der Eingabedatei angegeben ist (here)
	\item[t:] am oberen Ende der aktuellen oder Folgeseite (top)
	\item[b:] am unteren Ende der aktuellen Seite (bottom)
	\item[p:] auf einer eigenen Seite für ein oder mehrere Gleitobjekte (page)
	\item[!:] Überschreiben Sie die internen Parameter, die LaTeX zur Bestimmung "guter" Gleitkommapositionen verwendet.
	\item[H:] Setzt den Float an genau die Stelle im LaTeX-Code. Erfordert das float-Paket.
\end{description}

 %durch Eigene Kapitel erstezen
%\chapter{Einfügen von Quellcode}
\label{chap:code}

\section{Beispiel für einen Programmcode}


\if 1\mintedinstall
\subsection{Beispiel Minted}
	\begin{minted}[xleftmargin=20pt,linenos]{python}
	import numpy as np
	
	def incmatrix(genl1,genl2):
		m = len(genl1)
		n = len(genl2)
		M = None #to become the incidence matrix
		VT = np.zeros((n*m,1), int)  #dummy variable
	
		#compute the bitwise xor matrix
		M1 = bitxormatrix(genl1)
		M2 = np.triu(bitxormatrix(genl2),1) 
	
		for i in range(m-1):
			for j in range(i+1, m):
				[r,c] = np.where(M2 == M1[i,j])
				for k in range(len(r)):
					VT[(i)*n + r[k]] = 1;
					VT[(i)*n + c[k]] = 1;
					VT[(j)*n + r[k]] = 1;
					VT[(j)*n + c[k]] = 1;
					
					if M is None:
						M = np.copy(VT)
					else:
						M = np.concatenate((M, VT), 1)
					
					VT = np.zeros((n*m,1), int)
	
		return M
	\end{minted}
\fi




\subsection{Beispiel listings}
\begin{lstlisting}[language=Python]
import numpy as np

def incmatrix(genl1,genl2):
	m = len(genl1)
	n = len(genl2)
	M = None #to become the incidence matrix
	VT = np.zeros((n*m,1), int)  #dummy variable

	#compute the bitwise xor matrix
	M1 = bitxormatrix(genl1)
	M2 = np.triu(bitxormatrix(genl2),1) 
	
	for i in range(m-1):
		for j in range(i+1, m):
			[r,c] = np.where(M2 == M1[i,j])
			for k in range(len(r)):
				VT[(i)*n + r[k]] = 1;
				VT[(i)*n + c[k]] = 1;
				VT[(j)*n + r[k]] = 1;
				VT[(j)*n + c[k]] = 1;
	
				if M is None:
					M = np.copy(VT)
				else:
					M = np.concatenate((M, VT), 1)
				
				VT = np.zeros((n*m,1), int)
				
	return M
\end{lstlisting}


%\chapter{Einfügen von Tabellen}
\label{chap:tabellen}

\section{Beispieltabelle}

\begin{table}[h]
	\caption{Beispieltabelle}
	\label{tab:Parameter}
	\centering
	\begin{tabularx}{\textwidth}{p{0.32\textwidth}p{0.32\textwidth}p{0.32\textwidth}}
																					\toprule
		Eins                       	     & Zwei    			& Drei 				\\ 	\midrule
		Vier                       	     & Fünf    			& Sechs				\\
		Sieben                           & Acht         	& Neun				\\ 	\bottomrule	
	\end{tabularx}
\end{table}

\begin{longtabu} to \textwidth {@{} X X[c] X[r] @{}}
	\caption{Tabelle auf Textbreite mit drei gleich großen Spalten}\label{tab:bsp1} \\
	\toprule
	Spalte 1 linksbündig & Spalte 2 zentriert & Spalte 3 rechtsbündig \\
	\midrule
1 , 2                           & 3 , 2                           & 1 , 3 \\
2 , 4                           & 6 , 4                           & 2 , 6 \\
3 , 6                           & 9 , 6                           & 3 , 9 \\
	\bottomrule
\end{longtabu}

\begin{longtabu} to 140mm {@{} X X[c] X[r] @{}}
	\caption{Tabelle auf Textbreite mit drei gleich großen Spalten}\label{tab:bsp2} \\
	\toprule
	Spalte 1 linksbündig & Spalte 2 zentriert & Spalte 3 rechtsbündig \\
	\midrule
1 , 2                           & 3 , 2                           & 1 , 3 \\
2 , 4                           & 6 , 4                           & 2 , 6 \\
3 , 6                           & 9 , 6                           & 3 , 9 \\
	\bottomrule
\end{longtabu}
%

\begin{longtabu} to \textwidth {@{} X[1,l] X[1,c] X[1,r] @{}}
	%----- Kopfzeile erste Tabelle ----- %
	\caption{Tabelle über mehrere Seiten}\label{tab:mehrere Seiten} \\
	\toprule
	Spalte 1 linksbündig & Spalte 2 zentriert & Spalte 3 rechtsbündig \\
	\midrule
	\endfirsthead
	%----- Kopfzeile zweite Tabelle ----- %	
	\caption*{\textbf{Fortsetzung:} \cref{tab:mehrere Seiten}} \\
	\toprule
	Spalte 1 linksbündig & Spalte 2 zentriert & Spalte 3 rechtsbündig \\
	\midrule
	\endhead
	%----- Tabellenende erste Tabelle ----- %	
	\bottomrule
	\endfoot
	%----- Tabellenende zweite Tabelle ----- %
	\bottomrule    
	\endlastfoot
	%----- Inhalt der Tabelle Tabelle ----- %	
1 , 2                           & 3 , 2                           & 1 , 3 \\
2 , 4                           & 6 , 4                           & 2 , 6 \\
3 , 6                           & 9 , 6                           & 3 , 9 \\
4 , 8                           & 12 , 8                           & 4 , 12 \\
5 , 10                           & 15 , 10                           & 5 , 15 \\
6 , 12                           & 18 , 12                           & 6 , 18 \\
7 , 14                           & 21 , 14                           & 7 , 21 \\
8 , 16                           & 24 , 16                           & 8 , 24 \\
9 , 18                           & 27 , 18                           & 9 , 27 \\
10 , 20                           & 30 , 20                           & 10 , 30 \\
11 , 22                           & 33 , 22                           & 11 , 33 \\
12 , 24                           & 36 , 24                           & 12 , 36 \\
13 , 26                           & 39 , 26                           & 13 , 39 \\
14 , 28                           & 42 , 28                           & 14 , 42 \\
15 , 30                           & 45 , 30                           & 15 , 45 \\
16 , 32                           & 48 , 32                           & 16 , 48 \\
17 , 34                           & 51 , 34                           & 17 , 51 \\
18 , 36                           & 54 , 36                           & 18 , 54 \\
19 , 38                           & 57 , 38                           & 19 , 57 \\
\end{longtabu}
%\chapter{Mathematische Beispiele}
\label{sec:math}


\section{Gleichungen}

\begin{align}
\sin A \cos B &= \frac{1}{2}\left[ \sin(A-B)+\sin(A+B) \right] \\
\sin A \sin B &= \frac{1}{2}\left[ \sin(A-B)-\cos(A+B) \right] \\
\cos A \cos B &= \frac{1}{2}\left[ \cos(A-B)+\cos(A+B) \right] 
\end{align}

\begin{align*}
\sin A \cos B &= \frac{1}{2}\left[ \sin(A-B)+\sin(A+B) \right] \\
\sin A \sin B &= \frac{1}{2}\left[ \sin(A-B)-\cos(A+B) \right] \\
\cos A \cos B &= \frac{1}{2}\left[ \cos(A-B)+\cos(A+B) \right] 
\end{align*}

\[
\int_a^bu\frac{d^2v}{dx^2}\,dx
=\left.u\frac{dv}{dx}\right|_a^b
-\int_a^b\frac{du}{dx}\frac{dv}{dx}\,dx.
\]



\section{Arrays}

\[
\begin{bmatrix}
1 & x & 0 \\
0 & 1 & -1
\end{bmatrix}\begin{bmatrix}
1  \\
y  \\
1
\end{bmatrix}
=\begin{bmatrix}
1+xy  \\
y-1
\end{bmatrix}.
\]


\[
|x|=\begin{cases}
x, & \text{if }x\geq 0\,,  \\
-x, & \text{if }x< 0\,.
\end{cases}
\]


\[
\begin{matrix}
-2 & 1 & 0 & 0 & \cdots & 0  \\
1 & -2 & 1 & 0 & \cdots & 0  \\
0 & 1 & -2 & 1 & \cdots & 0  \\
0 & 0 & 1 & -2 & \ddots & \vdots \\
\vdots & \vdots & \vdots & \ddots & \ddots & 1  \\
0 & 0 & 0 & \cdots & 1 & -2
\end{matrix}
\]






%\chapter{tikz - Grafiken}
\label{sec:tikz}

\section{Beispielkapitel tikz - Grafiken}
\label{sec:beispiel}

\blindtext

\begin{figure}[!htbp]
\centering
\begin{align*}
f(x)=\tanh(x)
\end{align*}
\begin{tikzpicture}
\begin{axis}
[
width=200pt,
height=200pt,
axis x line=middle, xmin=-4, xmax=4, xtick={-3,...,3}, xlabel=$x$,
axis y line=middle, ymin=-1.5, ymax=1.5, ytick={-1,...,1}, ylabel=$f(x)$,
scale only axis=true,
samples=101
]
\addplot[blue,mark=none, very thick]{tanh(x)};
%		\legend{$\tanh(x)$}
\end{axis}
\end{tikzpicture}
\caption[Tangens hyperbolicus Aktivierungsfunktion]{Tangens hyperbolicus}
\label{fig:AktivTan}
\end{figure}

\blindtext


\begin{figure}[!htbp]
	\centering
	\begin{tikzpicture}[]
	\tikzstyle{netnode}=[circle, inner sep=0pt, text width=22pt, align=center, line width=1.0pt]
	\tikzstyle{inputnode}=[netnode, fill=lightgray,draw=black]
	\tikzstyle{hiddennode}=[netnode, fill=lightgray,draw=black]
	\tikzstyle{outputnode}=[netnode, fill=lightgray,draw=black]
	\tikzstyle{signal}=[arrows={-latex},draw=black,line width=1pt,rounded corners=4pt]
	
	\def\nodedist{35pt}
	\def\layerdist{80pt}
	\def\pindist{20pt}
	
	\tikzstyle{every pin edge}=[signal]
	\tikzstyle{annot} = [text width=6em, text centered]
	
	\foreach \y in {1,...,3}
	\node[inputnode, pin={[pin edge={latex-}, pin distance=\pindist]left:Eingabe \y}] 
	(I\y) at (0,-\y*\nodedist) {$i_\y$};  
	
	\foreach \y in {1,...,4}
	\node[hiddennode] 
	(H1\y) at ($(\layerdist,-\y*\nodedist) +(0, 0.5*\nodedist)$) {$h_\y^1$};
	
	\foreach \y in {1,...,4}
	\node[hiddennode] 
	(H2\y) at ($(2*\layerdist,-\y*\nodedist) +(0, 0.5*\nodedist)$) {$h_\y^2$};
	
	\foreach \y in {1,...,2}
	\node[outputnode, pin={[pin edge={-latex}, pin distance=\pindist]right:Ausgabe \y}]
	(O\y) at ($(H21) + (\layerdist, -\y*\nodedist)$) {$o_\y$};
	
	\foreach \dest in {1,...,4}
	\foreach \source in {1,...,3}
	\draw[signal] (I\source) -- (H1\dest);
	
	\foreach \dest in {1,...,4}
	\foreach \source in {1,...,4}
	\draw[signal] (H1\source) -- (H2\dest);
	
	\foreach \dest in {1,...,2}
	\foreach \source in {1,...,4}
	\draw[signal] (H2\source) edge (O\dest);
	
	\node[annot, above=4pt of H11] (hl) {verborgene Schicht 1};
	\node[annot, above=4pt of H21] (hl) {verborgene Schicht 2};
	\node[annot] at (I1 |- hl) {Eingabe\-schicht};
	\node[annot] at (O1 |- hl) {Ausgabe\-schicht};
	\end{tikzpicture}
	\caption{Schematischer Aufbau eines künstlichen neuronalen Netzes \cite[Abb. nach][]{Frochte.2019}}
	\label{fig:DNN}
\end{figure}


\section{Beispielkapitel Standard Grafik}
\label{sec:graf}

\blindtext

\begin{figure}[h]
  \centering  
  	\includegraphics[scale=0.5]{Images/ILEK-logo.jpg}
  \caption{ILEK Logo}
  \label{fig:starwars}
\end{figure}

\blindtext


%\begin{figure}[!htbp]
%	\begin{minipage}[c]{0.48\textwidth}
%		\input{Images/Bewertungsfunktionen/reward_photo_05_L00.pgf}
%		\centering\textsf{\textbf{(a)} Umgebungsbeleuchtung $\geq$ 500\,lx}
%	\end{minipage}
%	\hfill
%	\begin{minipage}[c]{0.48\textwidth}
%		\input{Images/Bewertungsfunktionen/reward_photo_05_L11.pgf}
%		\centering\textsf{\textbf{(b)} Keine Umgebungsbeleuchtung}
%	\end{minipage}
%	\caption{Bewertungsfunktion \gf{photo 05}}\label{fig:rewP05}
%\end{figure}


\begin{figure}
	\centering
	\begin{subfigure}[b]{0.3\textwidth}
		\centering
		\includegraphics[width=\textwidth]{Images/ILEK-logo.jpg}
		\caption{$y=x$}
		\label{fig:a}
	\end{subfigure}
%	\hfill
\hspace{1cm}
	\begin{subfigure}[b]{0.3\textwidth}
		\centering
		\includegraphics[width=\textwidth]{Images/ILEK-logo.jpg}
		\caption{$y=3sinx$}
		\label{fig:b}
	\end{subfigure}
	\\ 
	\begin{subfigure}[b]{0.3\textwidth}
		\centering
		\includegraphics[width=\textwidth]{Images/ILEK-logo.jpg}
		\caption{$y=3sinx$}
		\label{fig:c}
	\end{subfigure}
%	\hfill
\hspace{1cm}
	\begin{subfigure}[b]{0.3\textwidth}
		\centering
		\includegraphics[width=\textwidth]{Images/ILEK-logo.jpg}
		\caption{$y=5/x$}
		\label{fig:five over x}
	\end{subfigure}
	\caption{Three simple graphs}
	\label{fig:d}
\end{figure}






% ============= Anhang =============

\appendix
\begin{appendices}
	\chapter{Exemplarischer Anhang}
\label{chap:tab}

\section{Beispieltabelle}

\begin{table}[h]
	\caption{Beispieltabelle}
	\label{tab:par}
	\centering
	\begin{tabularx}{\textwidth}{p{0.32\textwidth}p{0.32\textwidth}p{0.32\textwidth}}
																					\toprule
		Eins                       	     & Zwei    			& Drei 				\\ 	\midrule
		Vier                       	     & Fünf    			& Sechs				\\
		Sieben                           & Acht         	& Neun				\\ 	\bottomrule	
	\end{tabularx}
\end{table}

\begin{longtabu} to \textwidth {@{} X X[c] X[r] @{}}
	\caption{Tabelle auf Textbreite mit drei gleich großen Spalten} \\
	\toprule
	Spalte 1 linksbündig & Spalte 2 zentriert & Spalte 3 rechtsbündig \\
	\midrule
1 , 2                           & 3 , 2                           & 1 , 3 \\
2 , 4                           & 6 , 4                           & 2 , 6 \\
3 , 6                           & 9 , 6                           & 3 , 9 \\
	\bottomrule
\end{longtabu}

\begin{longtabu} to 140mm {@{} X X[c] X[r] @{}}
	\caption{Tabelle auf Textbreite mit drei gleich großen Spalten} \\
	\toprule
	Spalte 1 linksbündig & Spalte 2 zentriert & Spalte 3 rechtsbündig \\
	\midrule
1 , 2                           & 3 , 2                           & 1 , 3 \\
2 , 4                           & 6 , 4                           & 2 , 6 \\
3 , 6                           & 9 , 6                           & 3 , 9 \\
	\bottomrule
\end{longtabu}

\begin{longtabu} to \textwidth {@{} X[1,l] X[1,c] X[1,r] @{}}
	%----- Kopfzeile erste Tabelle ----- %
	\caption{Tabelle über mehrere Seiten}\label{tab:mehrere Seiten Anhang} \\
	\toprule
	Spalte 1 linksbündig & Spalte 2 zentriert & Spalte 3 rechtsbündig \\
	\midrule
	\endfirsthead
	%----- Kopfzeile zweite Tabelle ----- %	
	\caption*{\textbf{Fortsetzung:} \Cref{tab:mehrere Seiten Anhang}} \\
	\toprule
	Spalte 1 linksbündig & Spalte 2 zentriert & Spalte 3 rechtsbündig \\
	\midrule
	\endhead
	%----- Tabellenende erste Tabelle ----- %	
	\bottomrule
	\endfoot
	%----- Tabellenende zweite Tabelle ----- %
	\bottomrule    
	\endlastfoot
	%----- Inhalt der Tabelle Tabelle ----- %	
	1 , 2                           & 3 , 2                           & 1 , 3 \\
	2 , 4                           & 6 , 4                           & 2 , 6 \\
	3 , 6                           & 9 , 6                           & 3 , 9 \\
	4 , 8                           & 12 , 8                           & 4 , 12 \\
	5 , 10                           & 15 , 10                           & 5 , 15 \\
	6 , 12                           & 18 , 12                           & 6 , 18 \\
	7 , 14                           & 21 , 14                           & 7 , 21 \\
	8 , 16                           & 24 , 16                           & 8 , 24 \\
	9 , 18                           & 27 , 18                           & 9 , 27 \\
	10 , 20                           & 30 , 20                           & 10 , 30 \\
	11 , 22                           & 33 , 22                           & 11 , 33 \\
	12 , 24                           & 36 , 24                           & 12 , 36 \\
	13 , 26                           & 39 , 26                           & 13 , 39 \\
	14 , 28                           & 42 , 28                           & 14 , 42 \\
	15 , 30                           & 45 , 30                           & 15 , 45 \\
	16 , 32                           & 48 , 32                           & 16 , 48 \\
	17 , 34                           & 51 , 34                           & 17 , 51 \\
	18 , 36                           & 54 , 36                           & 18 , 54 \\
	19 , 38                           & 57 , 38                           & 19 , 57 \\
\end{longtabu}

\end{appendices}


% ============= Verzeichnise =============

\printbibliography[title=Literaturverzeichnis] %Literaturverzeichnis
\addcontentsline{toc}{chapter}{Literaturverzeichnis}

\listoffigures %Verzeichnis aller Bilder
\addcontentsline{toc}{chapter}{Abbildungsverzeichnis}

\listoftables %Verzeichnis aller Tabellen
\addcontentsline{toc}{chapter}{Tabellenverzeichnis}


% ============= Dokumentende =============

\end{document}
